\scanpage{101}
\mspage{49}
% Main text, approximately lines 1--2.
En résumé, on a un foncteur \deleted{\illegible{1}}\inserted{canonique}
% Approximately lines 3--10; the source of the vertical functor.
Cat. isotopique des \((X,S,\Gamma)\), \(X\) surface compacte \deleted{\illegible{1}},
\unsure{\(\Gamma\)} orienté, \(S\) partie discrète, \(\Gamma\) opérant par \unsure{\(\chi\)-autom.}
(\inserted{\(\sigma\in\Gamma\)}\unsure{respectant|respecte} l'orientation si \(\chi(\sigma)=+1\), la renversant
si \(\chi(g)=-1\)), \(\Gamma\) transitif sur \(\Pi_0(X)\),
\deleted{\illegible{1}} et si \((X_0,S_0)\) est une composante connexe
de \((X,S)\), on veut que \deleted{\((X_0,S_0)\) \illegible{2}}\inserted{si \(X_0\simeq S^2\), on}
\unsure{\uline{ait}|\deleted{ait}} \(\operatorname{card} S_{\unsure{0}}\geq 3\) \deleted{\illegible{3}}.
\margin{flèches : \unsure{classes} d'isotopie d'homéom. \unsure{respectant} toutes les structures.}

\diagram{\begin{tikzcd}[row sep=2.5em]
{} \arrow[d] \\
{}
\end{tikzcd}}

% Approximately lines 11--18.
Catégorie des groupes extérieurs à \inserted{\(\Gamma\)-}lacets \(\pi\),
\inserted{\(\pi\) non commutatif (i.e. \(\pi\neq 0,\mathbb Z\));}
sur lesquels un sous-groupe \(\Sigma\subset\Gamma\) de \(\Gamma\) \inserted{opère}
(cette description étant
équivalente : une description en termes
de systèmes de groupoïdes
\[
\Pi_{\unsure{D^*},\Gamma}\longrightarrow\Pi_{U,\Gamma}
\longrightarrow\Pi_{e,\Gamma}\quad\text{et }\kappa\ldots,
\]
plus conceptuelle \unsure{des contraintes} \illegible{1}).
\margin{Attention \unsure{au} cas \(S=\varnothing\) dans la définition des structures : la donnée de
\(H^2(\pi,\mathbb Z)\overset{\kappa}{\simeq}T\)); il faut bien sûr tenir compte \unsure{lors de} l'action de \(\Sigma\) sur
\(H^2(\pi,\mathbb Z)\) \deleted{\unsure{et sur \(T\)}}
(via son action ext. sur \(\pi\)) et sur \(T\) (via \(\chi\)) \illegible{2} compatibilité.}
\margin{\unsure{Ça revient d'ailleurs} \unsure{au cas de} \unsure{foncteurs} \illegible{1}.}

% Approximately lines 19--25; the sentence continues on scan 102.
\deleted{\unsure{Pour déf.}} Je présume que la démonstration
\unsure{du} fait que ce foncteur soit pl. fidèle
ne fera pas de difficultés essentielles,
en utilisant ce qui est connu pour
\(\Gamma=1\). Mais le fait que, pour \(\Gamma\) groupe
fini donné, \deleted{\unsure{qu'il soit un isomorp.}}\inserted{il soit \unsure{ess\textsuperscript{t}} surjectif}
% Lower-left handwritten pagination; attribution of this numbering is uncertain.
\margin{100}

\scanpage{102}
% Approximately lines 1--7; continuation from scan 101.
est un problème ouvert sur lequel les
gens sèchent. Bien sûr OPS \(\Sigma=\Gamma\),
\inserted{et \(\Gamma\subset\) Teichmüller étendu}
et la question est si tt sous-groupe
fini du Teichmüller étendu se
réalise comme groupe opérant sur \((X_0,S_0)\) \illegible{1}
de façon essentiellement unique. Plus
précisément, si \deleted{\illegible{1}} \(A=\) groupe des \deleted{\unsure{diffé}}homéomorphismes
% Approximately lines 8--16.
(ou difféomorphismes, si \(X_0\) est
différentiable) de \(X_0\), \(A^0\) sa compos. \unsure{de}
connexe (\uline{NB} \(A^0\) est contractile, \deleted{\unsure{)}}\inserted{dans}
les cas anabéliens) donc \(A/A^0=\mathcal T_{g,\nu}\)
(groupe de Teichmüller pour genre \(g\) et \(\nu\)
trous), la question revient : \unsure{ceci} si pour tt
hom. d'un groupe fini \(\Sigma\) dans \deleted{\illegible{1}} \(\mathcal T_{g,\nu}\),
(OPS \(\Sigma\subset\mathcal T_{g,\nu}\)), celui-ci se relève
en un hom. dans \(A\), et si deux tels
relèvements sont conjugués par un \(a\in A^0\)
(isotopie \inserted{au sens strict} de deux relèvements \unsure{\(\sim\)}).

% Approximately lines 17--21; continuation on scan 103.
Ayant abouti à une \deleted{description}\inserted{\unsure{interprétation}} tellement
simple de la 1-catégorie (``1-isotopique'') déduite
de la 2-catégorie des systèmes
\((\Pi_{D^*,\Gamma}\longrightarrow\Pi_{U,\Gamma}\longrightarrow\Pi_{e,\Gamma},\kappa)\), en termes
% Lower-left additional handwritten pagination; hand/status uncertain.
\margin{101}

\scanpage{103}
\mspage{50}
% Approximately lines 1--7; continuation from scan 102.
de groupes extérieurs à lacets (\inserted{\(\pi\)}munis
d'un ens de hom. ext. \(\kappa_i:T\longrightarrow\pi\), et
à défaut d'un \(\kappa:T\longrightarrow H^2(\pi,\mathbb Z)\))
(pour bien faire, il faudrait écrire
\(T^{(\otimes-1)}\simeq H^2(\pi,\mathbb Z)\), mais ici on a un iso.
canonique \(T^{\otimes-1}\simeq T\) i.e. \(T^{\otimes 2}\simeq\mathbb Z\ldots\))],
% Approximately lines 8--14.
la question se pose \unsure{de comment récupérer},
(à équivalence définie à isom unique près)
ce diagramme, en termes de \(\kappa\); tout
revient à la \deleted{définition} description des
catégories \(\Pi_{D^*,\unsure{\Gamma}}\) et \(\Pi_{U,\Gamma}\), et des deux
foncteurs \(\Pi_{D^*,\Gamma}\longrightarrow\Pi_{U,\Gamma}\), \(\Pi_{U,\Gamma}\longrightarrow\Pi_{e,\Gamma}\);
\inserted{ou ce qui revient au \(=\), des topos (multigaloisiens) et morphismes de topos correspondants.}

% Approximately line 15.
a) \uline{Description} de \(\Pi_{U,\Gamma}\) et de \(\Pi_{U,\Gamma}\longrightarrow\Pi_{e,\Gamma}\).

% Approximately lines 16--22; continuation on scan 104.
Soit plus généralement \(\pi\) un groupe \inserted{extérieur}, \deleted{\unsure{tel que}} \unsure{dont le}
\uline{centre} \uline{soit trivial} (ceci correspond à l'hyp.
anabélienne !), montrons comment on
lui associe un topos classifiant
\(B_\pi\), qui (comme cat. de faisceaux) sera \inserted{\unsure{pour} \illegible{1}}
\(\operatorname{Ens}(\pi)\), de façon ``fonctorielle''. Tout
revient à voir comment, à une \uline{classe}
% Lower-left additional handwritten pagination; hand/status uncertain.
\margin{102}

\scanpage{104}
% Approximately lines 1--7; continuation from scan 103.
de \uline{conjugaison} \inserted{\(\Theta\)} d'iso
\[
\pi'\xrightarrow{u}\pi,
\]
on associe un foncteur ``image inverse''
\(\operatorname{Ens}(\pi)\longrightarrow\operatorname{Ens}(\pi')\), défini à isom unique
près. Considérons pour tt \(u\) de la classe
\(\Theta\), le foncteur ``\(u\)-restriction des op.''
\[
\operatorname{Ens}(\pi)\xrightarrow{u^*}\operatorname{Ens}(\pi')
\]
% Approximately lines 8--11, including interlinear corrections.
[qui \deleted{\unsure{est}}\inserted{définit, avec \(u_*=(u^{-1})^*\),}
\unsure{donc} (\unsure{une} \deleted{\unsure{isomorphisme}}\inserted{équivalence} des Topos
\[
B_{\pi'}\xrightarrow[\sim]{\unsure{u_\bullet|u_*}} B_\pi
\]
]
On va, entre ces \deleted{foncteurs}\inserted{équivalences}, \inserted{pour \(u\in\Theta\),} définir un
syst. transitif d'iso (ce qui permet
de les identifier entre eux !).

% Approximately lines 12--20.
\unsure{Soient|Soit} \(u,u'\in\Theta\), d'où \(u^*,u'^*\); on a par hyp.
\deleted{définition} un \(g\in\pi\) tel que \(u'=\operatorname{int}(g)\circ u\),
de plus \(g\) est déterminé mod un
élément de \(\operatorname{Centr}_\pi u(\pi')=\operatorname{Centr}(\pi)=1\);
donc ici \(g\) est unique. Mais \(g\)
peut servir à définir un isom
fonctoriel \deleted{entre \(u^*\) et \(u'^*\)}\inserted{entre}
\[
i_{u',u}:u^*\xrightarrow{\sim}u'^*
\]
en prenant, pour \(E\in\operatorname{Ob}\operatorname{Ens}(\pi)\),
\[
i_{u',u}(E):
\underset{\substack{\unsure{\Vert}\\E}}{u^*(E)}
\longrightarrow
\underset{\substack{\unsure{\Vert}\\E}}{u'^*(E)},
\qquad i_{u',u}(E)=g_E.
\]
% Lower-left additional handwritten pagination; hand/status uncertain.
\margin{103}

\scanpage{105}
\mspage{51}
% Approximately lines 1--6.
\uline{NB} Le raisonnement marche pour tte classe
de conj. d'hom. \inserted{de groupes} \(\pi'\longrightarrow\pi\) (i.e. tt hom.
ext.) dont le centralisateur dans \(\pi\) est
réduit à 1 (ce qui pour une \unsure{iso} se
\unsure{réduit à l'}hyp \(\operatorname{Centr}(\pi)=1\)).
\margin{\unsure{marche pour} les \deleted{\unsure{applications}} \unsure{hom. de tels \(\pi\)}
dont l'image soit d'indice fini (car le centralisateur dans un tel groupe \(\pi\)
d'un sous-groupe \unsure{d'i-} fini est encore \unsure{réduit à 1}).}

% Approximately lines 7--14.
Si maintenant un groupe \(\Sigma\) opère
sur le groupe extérieur \inserted{\(\pi\)}, alors par le
résultat précédent on peut dire qu'il
opère aussi sur le topos \(B_\pi\), \inserted{d'où un topos \(B_{\pi,\unsure{\Gamma|\Sigma}}\),} (on peut
dire aussi que \deleted{\illegible{1}} Comme
\(\operatorname{Centr}(\pi)=1\)) l'opération de \deleted{\(\unsure{\pi}\)}\(\Sigma\) sur \(\pi\)
définit une ext. \(\mathcal E\) de \(\Sigma\) par \(\pi\),
d'où un topos \(B_{\pi,\unsure{\Gamma|\Sigma}}=B_{\mathcal E}\), \deleted{\unsure{et alors}}\inserted{\unsure{On récupère}}
d'où aussi \(B_{\pi,\unsure{\Gamma|\Sigma}}\longrightarrow B_{\unsure{\Gamma|\Sigma}}\). Pour se
% Approximately lines 15--22; continuation on scan 106.
tranquilliser, il faudrait s'assurer que
si on a un hom de groupes extérieurs
\(\pi'\xrightarrow{\Theta}\pi\) commutant à l'action d'un
\(\Sigma\), avec \(\operatorname{Centr}_\pi(\Theta)=(1)\), alors il existe
un \(\Theta_{\unsure{\Gamma|\Sigma}}:B_{\pi',\unsure{\Gamma|\Sigma}}\longrightarrow B_{\pi,\unsure{\Gamma|\Sigma}}\) défini à isom
canonique près. Or tt \(u\in\Theta\) définit un
hom d'extensions \inserted{\(u_{\unsure{\Gamma|\Sigma}}\) (de \(\Sigma\) par \(\pi'\) resp. \(\pi\))}
\(\mathcal E'\longrightarrow\mathcal E\) en dessous de
\(u\), et si on passe de \(u\) à \(u'=\operatorname{int}(g)\circ u\), on aura
% Lower-left additional handwritten pagination; hand/status uncertain.
\margin{104}

\scanpage{106}
% Approximately lines 1--2; continuation from scan 105.
\(u'_{\unsure{\Gamma|\Sigma}}=\operatorname{int}_{\mathcal E}(g)\circ u_{\unsure{\Gamma|\Sigma}}\), et on termine comme
plus haut, avec unicité de \(g\).

% Approximately lines 3--8.
Soit maintenant un hom. extérieur \inserted{\uline{injectif}}
\[
\deleted{\unsure{\lambda}_{\unsure{\kappa}}}\Lambda=[\kappa]:T\longrightarrow\pi
\qquad (T\text{ commutatif}),
\]
supposons (pour \(\kappa\in[\kappa]\))
\[
\operatorname{Centr}_\pi\kappa=\kappa(T)
\]
(ce qui ne dépend pas du choix de \(\kappa\) dans
\([\kappa]\)). Je vais alors définir une
\margin{L'injectivité dans le cas qui nous intéresse résulte de l'hyp. anabélienne
(\unsure{infinité des ``lacets''}).}
% Approximately lines 9--17.
\deleted{catégorie} groupoïde abélien connexe \inserted{\(\Pi_\Lambda\)}, et un
\unsure{iso} de son \(\pi_{\unsure{1}}\) avec \(T\) \deleted{Un objet}
(d'où un topos, qui joue le rôle de
\(\widehat{\Pi}_{\unsure{D|D^*}}\)). Un objet sera un \(\kappa\in[\kappa]\).
Un hom de \(\kappa\) en \(\kappa'\) sera un
élément \(g\in\operatorname{Transp}_\pi(\kappa,\kappa')\). \deleted{Soit \(g\)}
La composition des hom évidente. On
trouve bien un groupoïde connexe, dans
lequel \(\operatorname{Aut}(\kappa)=\operatorname{Centr}_\pi(\kappa)\overset{\inserted{\text{par hyp.}}}{=}\kappa(T)\unsure{\xleftarrow[\sim]{\kappa}}T\).

% Approximately lines 18--21; continuation on scan 107.
OK. Bien sûr, on a un hom
\(\Pi_\Lambda\longrightarrow\deleted{\unsure{\pi}}\) groupoïde ponctuel \inserted{défini par \(\pi\)}, d'où
un hom topos \unsure{classifiant|classifiants} \inserted{\(B_\Lambda\longrightarrow B_\pi\)}.
Il faut voir la compatibilité de cette
% Lower-left additional handwritten pagination; hand/status uncertain.
\margin{105}

\scanpage{107}
\mspage{52}
% Approximately lines 1--7; continuation from scan 106.
construction par hom. extérieurs. Soit donc
\[
\Theta=f_\pi\deleted{\unsure{\Theta}}=[u]:\pi'\longrightarrow\pi
\]
un hom extérieur tel que \(\operatorname{Centr}_\pi(\Theta)=1\),
et \(\Lambda':T\longrightarrow\pi'\) un hom extérieur injectif
tel que \(\kappa'\in\Lambda'\Rightarrow\operatorname{Centr}_{\pi'}(\kappa')=\operatorname{Im}\kappa'\)
et considérons le composé
\[
\deleted{\unsure{\Lambda=}}f_\pi\circ\Lambda':T\longrightarrow\pi.
\]
% Approximately lines 8--12.
Soit \(d\in\mathbb Z\) et considérons
\[
\Lambda=f_\pi\circ\Lambda'\circ(d.\operatorname{id}_T)
\]
\inserted{supposons que \(\kappa\in\Lambda\Rightarrow\operatorname{Centr}_\pi\kappa=\operatorname{Im}\kappa\)}
On va définir un hom de \(\Pi_{\Lambda'}\) dans \(\Pi_\Lambda\),
d'où un hom de Topos \(B_{\Lambda'}\longrightarrow B_\Lambda\),
et une donnée de commutation \inserted{\(\alpha\)} du
\diagram{\begin{tikzcd}[ampersand replacement=\&,column sep=4em,row sep=3.8em]
\& B_{\Lambda'} \arrow[r,"f_D"] \arrow[d] \& B_\Lambda \arrow[d] \\
(*) \& \unsure{(B,\pi')} \arrow[r,"{\begin{gathered}\deleted{\unsure{B(f_\pi)}}\\[-2pt]\inserted{B(f_\pi)\ \text{défini plus haut}}\end{gathered}}"']
\arrow[ur,bend left=32,shorten <=2.1em,shorten >=2.1em,"\alpha"']
\& \unsure{(B,\pi)}
\end{tikzcd}}
\margin{\illegible{2} \unsure{évidemment} \illegible{5} \deleted{\unsure{de Topos}} \illegible{2}.}
\deleted{\illegible{6} \(\Lambda',\Theta,d\)}

% Approximately lines 13--19; continuation on scan 108.
\deleted{Soit} \(\kappa':T\longrightarrow\pi'\) un objet de \(\Pi_{\Lambda'}\),
on veut définir (à isom unique près) un
objet \(\kappa\) de \(\Pi_\Lambda\). Pour tt \(u\in f_\pi=\Theta\),
on considère \(u\circ\unsure{\kappa|\kappa'}\circ(d.\operatorname{id}_T)\), il y a entre eux
un système transitif d'iso pour \(u\) variable
dans \(\Theta\), on peut les identifier entre eux.
% Lower-left additional handwritten pagination; hand/status uncertain.
\margin{106}

\scanpage{108}
% Approximately lines 1--7.
La fonctorialité de cet objet p.r. à
\(\kappa'\) variable est évidente : [on peut dire que
\(u\in\Theta\) définit un hom de groupoïdes
\[
u_\Lambda:\Pi_\Lambda\longrightarrow\Pi_{\Lambda'},
\]
et entre ceux-ci
il y a un syst. transitif d'iso \unsure{\}}. En fait,
pour \(u\) fixe on a un \deleted{syst.} diagramme
commutatif d'hom de groupoïdes
\diagram{\begin{tikzcd}[ampersand replacement=\&]
\Pi_{\Lambda'} \arrow[r,"u_\Lambda"] \arrow[d] \& \Pi_{\unsure{\Lambda|\Lambda e}} \arrow[d] \\
(e,\pi') \arrow[r,"u"] \& (e,\pi)
\end{tikzcd}}
% Approximately lines 8--11.
\deleted{D'où une}
et entre ces diagrammes il y a un
syst. transitif d'isomorphismes, d'où
le diagramme \((*)\) et la donnée
de commutation \(\alpha\).

% Approximately lines 12--19; the final note continues on scan 109.
\deleted{Le lien avec}
Quand on a une \inserted{famille \(\Lambda'\) (\unsure{ou un ens})} \deleted{\illegible{1}} d'hom
extérieurs \(\Lambda'_{i'}\ (i'\in I'):T\longrightarrow\pi'\), et une
famille \(\Lambda\unsure{=}(\Lambda_i)_{i\in I}\) d'hom ext. \(T\longrightarrow\pi\), et des
\(\tau:I'\longrightarrow I\), \deleted{\illegible{1}}\inserted{et \((d_{i'})_{i'\in I'}\)}
\inserted{\unsure{des \(d_{i'}\in\mathbb Z\), tels que}}
(pour \(\Theta:\pi'\longrightarrow\pi\) hom ext. donné)
\(\forall i'\in I'\), posant \(i=\tau(i')\), on ait
\[
\Lambda_i=\unsure{\Lambda'_{i'}|\Lambda_{i'}}\circ\Theta\circ(d_{i'}.\operatorname{id}).
\]
[\uline{NB} si \(I\longrightarrow\operatorname{Homext}(T,\pi)\) injectif, ces conditions
\unsure{impliquent} que \(\tau\) est déterminé par \(\Theta\), et si
% Lower-left additional handwritten pagination; hand/status uncertain.
\margin{107}

\scanpage{109}
\mspage{53}
% Approximately lines 1--6; continuation from scan 108.
de plus \(\mathbb Z\) opère fidèlement dans \(T\), p.ex.
si \(T\simeq\mathbb Z\), alors \((d_{i'})\) est également unique \(\ldots\)]
Alors on construit \deleted{\unsure{B}}
\[
\Pi_{D^*,\Lambda'}=
\deleted{\text{\unsure{Produit}}}\inserted{\text{somme}} \text{des }\inserted{\text{groupoïdes }}\Pi_{\Lambda'_{i'}}
\]
et de \(=\) \(\Pi_{D^*,\Lambda}\), et on trouve un
diagramme \inserted{iso}commutatif de topos
\diagram{\begin{tikzcd}[ampersand replacement=\&]
B_{D^*,\Lambda'} \arrow[r] \arrow[d] \& B_{D^*,\Lambda} \arrow[d] \\
B_{\pi'} \arrow[r] \& B_\pi
\end{tikzcd}}

% Approximately lines 7--10.
Supposons maintenant (enfin !) que
\(\Sigma\) opère sur ces groupes extérieurs à lacets \(\ldots\)
Je déclare forfait --- il est évident que
tout marche bien !

% Approximately lines 11--17.
B. \uline{Retour au cas arithmétique}, si on
veut décrire en termes ``galoisiens'' les
couples \((X,S)\) \deleted{\unsure{conn.}} anabéliens \inserted{connexes} sur un corps
\(K\) de t.f. sur \(\mathbb Q\). [\uline{NB} si on prend
un \(K\) de t.f. sur \(\mathbb F_p\), il faudrait se
borner aux groupes fondamentaux ``premiers à
\(p\)'', à \unsure{cel'|cela} près \unsure{un développement pourrait|les développements pourraient}
se \uline{faire} quand \(=\) \(\ldots\)]
% Lower-left additional handwritten pagination; hand/status uncertain.
\margin{108}

\scanpage{110}
% Approximately lines 1--7.
Il est donc clair qu'en termes d'une
clôture algébrique \(\overline K\) de \(K\), et du
groupe de galois \inserted{profini} \(\Gamma=\operatorname{Gal}(\overline K/\unsure{K|k})\), la
description la plus simple est en
termes de groupes extérieurs à lacets, et
d'actions extérieures de sous-groupes \inserted{ouverts} \(\Sigma\) de \(\Gamma\)
donnés. De façon précise, \deleted{soit} on
% Approximately lines 8--14.
choisit une composante connexe
\(\overline X_0\) de \(\overline X\) (ou ce qui revient au \(=\),
\(\overline U_0\) de \(\overline U=(X\setminus S)_{\overline K}\)), soit \(\Sigma\) son
stabilisateur dans \(\Gamma\) (il est remplacé
par un conjugué, quand on change
\(\overline U_0\)). Alors \(\Sigma\) opère sur le schéma
\(\unsure{\overline U_0|\overline X_0}\), donc opère extérieurement sur son
\(\pi_1\) (considéré comme groupe extérieur).
% Approximately lines 15--18.
Or sur celui-ci il y a une \inserted{\(T(\overline K)\)-}structure
à lacets, \deleted{définie par les}\inserted{\unsure{comme}} \unsure{un ens} d'indices
\[
I=S_0(\overline K)\subset \begin{array}[t]{@{}c@{}}S(\overline K)\\ \rotatebox{90}{$\simeq$}\\ S_0(\overline K)\wedge_\Sigma\Gamma\end{array}
\qquad[\text{vide \unsure{si|ssi} }S=\varnothing]
\]
% Lower-left additional handwritten pagination; hand/status uncertain.
\margin{109}
